\chapter{Introduction}

Few concept in physics have been proven as fruitful as Noether's theorem and the profound connection it revealed between symmetries and physical laws \cite{Noether1918}. It is no exaggeration to say that a large part of modern physics, both from experimental and theoretical standpoint, is devoted either to testing fundamental symmetries or to working out their various consequences in order to describe the multitude of phenomena occurring in Nature, from the microscopic to the cosmological scale. From the birth to the quantum mechanics, during a century of continuous efforts both in theoretical and experimental physic, the Standard Model has emerged as our most advanced theory for describing nature. One of its essential pillars is precisely this legacy: the various symmetries, and their related conservation rules that constituted the mathematical framework to describe the four fundamental forces of nature.
Although successful over the course of the years and able to describe with great precision many of the phenomena, we know that Standard Model is not definitive, and certain phenomena are either not fully explained by the current theory (for instance, the problem of the mass of neutrinos) or it is unable to explain the baryogenesis and the matter antimatter unbalance \cite{Dine_2003}. The two cases cited do not exhaust the tensions observed between theory and experimental observation, nor the problem of the internal consistency that theoretical physicist are trying to resolve. Nevertheless the matter antimatter unbalance represents an open question of enormous importance. Years of experiments at accelerator and cosmological observation have yet to provide a valid mechanism which can explain the matter antimatter unbalance.
Historically, the first coherent and systemic study of the matter-antimatter puzzle has been formulated by the russian physicist Andrei Sakharov in 1967 \cite{Sakharov1967}. He initially proposed that the observed baryon asymmetry could have arisen dynamically from an initially symmetric universe, provided three conditions were simultaneously satisfied:

\begin{itemize}
\item baryon number violation.
\item charge conjugation symmetry, C, and charge-parity conjugation symmetry CP violated simultaneously.
\item departure from thermal equilibrium.
\end{itemize}


\section{CPT symmetry and its implication}